% !TeX spellcheck = pt_BR
\documentclass[12pt,portuguese]{article}

%----------------------------------------------------------------------
% PACOTES ESSENCIAIS E TRADUÇÃO
%----------------------------------------------------------------------
\usepackage[brazilian]{babel}     % Suporte à língua portuguesa
\usepackage[utf8]{inputenc}       % Codificação de entrada
\usepackage[T1]{fontenc}          % Codificação da fonte
\usepackage{lmodern}              % Fonte escalável
\usepackage{microtype}            % Microtipografia (expansão e protrusão)

%----------------------------------------------------------------------
% PACOTES MATEMÁTICOS
%----------------------------------------------------------------------
\usepackage{amsmath, amsthm, amssymb, amstext, mathtools}
\usepackage{mathrsfs}             % Para fontes como \mathscr
\usepackage{dsfont}               % Para o indicador de conjunto \mathds{1}
\usepackage[all]{xy}              % Para diagramas comutativos
\usepackage{imakeidx}             % Criação de Lista de Símbolos / Índice

% Desigualdades sempre com símbolos inclinados
\let\leq\leqslant
\let\le\leqslant
\let\ge\geqslant
\let\geq\geqslant

%----------------------------------------------------------------------
% FORMATAÇÃO, LAYOUT E CORES
%----------------------------------------------------------------------
\usepackage{geometry}
\usepackage{varwidth}
\usepackage{graphicx}
\usepackage[pdftex]{xcolor}
\usepackage{titlesec}
\usepackage{appendix}
\usepackage{enumitem}
\usepackage{float}
\usepackage{booktabs}
\usepackage{array}
\usepackage{colortbl}
\usepackage{fontawesome5}
\usepackage{lucide-icons}
\usepackage{caption}
\usepackage{stackengine}
\captionsetup{width=0.8\textwidth}

% Definição de Cores
\definecolor{teal}{RGB}{15,118,110}
\definecolor{text}{RGB}{28,25,23}
\definecolor{muted}{RGB}{87,83,78}
\definecolor{card}{RGB}{255,255,255}
\definecolor{border}{RGB}{231,229,228}

% Definição de margens e espaçamento
\setlength{\topmargin}{-0.7in}
\setlength{\oddsidemargin}{0in}
\setlength{\textheight}{9.55in}
\setlength{\textwidth}{6.5in}
\setlength{\baselineskip}{12mm}

% Formatação dos títulos de seção
\titleformat{\section}{\Large\bfseries}{\thesection.}{0.5em}{}
\titleformat{\subsection}{\large\bfseries}{\thesubsection.}{0.5em}{}
\titleformat{\subsubsection}{\normalfont\itshape}{\thesubsubsection.}{0.6em}{}

%----------------------------------------------------------------------
% TIKZ E GRÁFICOS
%----------------------------------------------------------------------
\usepackage{tikz}
\usepackage{pgfplots}
\pgfplotsset{compat=1.17}
\usetikzlibrary{positioning,calc,fit,backgrounds,nfold,decorations.pathreplacing,calligraphy,tikzmark,patterns,arrows.meta,bending,hobby}

%----------------------------------------------------------------------
% REFERÊNCIAS, LINKS E METADADOS DO PDF
%----------------------------------------------------------------------
\usepackage[pdftex, colorlinks, linkcolor={red!50!black}, 
citecolor={red}, urlcolor={blue!80!black}]{hyperref}

\hypersetup{
pdftitle={O Processo de Poisson},
pdfauthor={L. Cioletti},
pdfsubject={processos estocásticos a tempo contínuo, processo de Poisson, distribuição exponencial, falta de memória, tempos entre chegadas},
pdfkeywords={processo de Poisson, processos estocásticos, tempo contínuo, distribuição exponencial, falta de memória, tempos entre chegadas, MAT-UnB}
}

% --- PACOTE E COMANDO PARA A LICENÇA CREATIVE COMMONS ---
\usepackage[type={CC},modifier={by-nc-sa},version={4.0}]{doclicense}
\newcommand\blfootnote[1]{%
\begingroup
\renewcommand\thefootnote{}\footnote{#1}%
\addtocounter{footnote}{-1}%
\endgroup
}

%----------------------------------------------------------------------
% ALIASCNT E AMBIENTES DE TEOREMA (Cleveref)
%----------------------------------------------------------------------
\usepackage{aliascnt}
\theoremstyle{definition}
\newtheorem{theorem}{Teorema}

\newaliascnt{definition}{theorem}
\newtheorem{definition}[definition]{Definição}
\aliascntresetthe{definition}

\newaliascnt{example}{theorem}
\newtheorem{example}[example]{Exemplo}
\aliascntresetthe{example}

\newaliascnt{proposition}{theorem}
\newtheorem{proposition}[proposition]{Proposição}
\aliascntresetthe{proposition}

\newaliascnt{corollary}{theorem}
\newtheorem{corollary}[corollary]{Corolário}
\aliascntresetthe{corollary}

\newaliascnt{lemma}{theorem}
\newtheorem{lemma}[lemma]{Lema}
\aliascntresetthe{lemma}

\newaliascnt{remark}{theorem}
\newtheorem{remark}[remark]{Observação}
\aliascntresetthe{remark}

% O primeiro argumento é o espaço acima do ambiente. Acrescentamos 0,4 cm ao
% valor usual de 3pt para separar melhor os exercícios.
\newtheoremstyle{plainexercise}
  {14.4pt}{3pt}{\normalfont}{}{\bfseries}{.}{.5em}
  {\thmnumber{#2}\thmnote{ #3}}
\theoremstyle{plainexercise}
\newtheorem{exercise}{Exercício}

% Versão marcada dos exercícios: a estrela fica pendurada na margem esquerda,
% como nas notas anteriores.
\newtheoremstyle{starredexercise}
  {14.4pt}{3pt}{\normalfont}{}{\bfseries}{.}{.5em}
  {\llap{\textasteriskcentered\,}\thmnumber{#2}\thmnote{ #3}}
\theoremstyle{starredexercise}
\newtheorem{starredexercise}[exercise]{Exercício}
\theoremstyle{definition}

% Ambientes não numerados usados no texto
\newtheorem*{property}{Propriedade}
\newtheorem*{solution}{Solução}

% Se um enumerate for o primeiro conteúdo de um teorema/exercício (cabeçalho
% ainda pendente, \@inlabel true), o \leavevmode dispara o \everypar do amsthm
% que coloca o cabeçalho na própria linha, antes da lista iniciar.
\makeatletter
\setlist[enumerate]{before={\if@inlabel\leavevmode\fi}}
\makeatother

% Nomes em português para os ambientes
\AtBeginDocument{\renewcommand{\proofname}{\normalfont\textbf{Prova}}}
\renewcommand{\qedsymbol}{$\blacksquare$}

% Cleveref (deve vir depois de hyperref e amsthm)
\usepackage[brazilian,nameinlink]{cleveref}

\crefname{theorem}{teorema}{teoremas}
\Crefname{theorem}{Teorema}{Teoremas}
\crefname{definition}{definição}{definições}
\Crefname{definition}{Definição}{Definições}
\crefname{lemma}{lema}{lemas}
\Crefname{lemma}{Lema}{Lemas}
\crefname{proposition}{proposição}{proposições}
\Crefname{proposition}{Proposição}{Proposições}
\crefname{corollary}{corolário}{corolários}
\Crefname{corollary}{Corolário}{Corolários}
\crefname{example}{exemplo}{exemplos}
\Crefname{example}{Exemplo}{Exemplos}
\crefname{exercise}{exercício}{exercícios}
\Crefname{exercise}{Exercício}{Exercícios}
\crefname{remark}{observação}{observações}
\Crefname{remark}{Observação}{Observações}
\crefname{section}{seção}{seções}
\Crefname{section}{Seção}{Seções}

%----------------------------------------------------------------------
% MACROS ADICIONAIS
%----------------------------------------------------------------------
\let\oldemptyset\emptyset
\let\emptyset\varnothing
\graphicspath{{./}{../arquivos-aux/}}

% Seno na convenção em português
\renewcommand{\sin}{\operatorname{sen}}

% Evita hifenização e quebra de sílabas no abstract mantendo a justificação
\expandafter\def\expandafter\abstract\expandafter{%
  \abstract
  \hyphenpenalty=10000
  \exhyphenpenalty=10000
  \emergencystretch=1.5em
}

%----------------------------------------------------------------------
% DOCUMENTO PRINCIPAL
%----------------------------------------------------------------------
\begin{document}

% --- CABEÇALHO COM O LOGO DA UnB ---
\vspace*{-0.8cm}
\begin{figure}[H]
\includegraphics{logo-mat-unb.pdf}
\end{figure}

% --- TÍTULO, AUTOR E DATA ---
\vspace*{0.3cm}
\begin{center}
{\LARGE\bf O Processo de Poisson}
\vspace{1cm}

{\large \bf L. Cioletti}
\vspace{0.3cm}

{\footnotesize Novembro de 2026}
\end{center}

% --- LICENÇA NO RODAPÉ ---
\blfootnote{%
\doclicenseImage[imagewidth=2.6em]\
\begingroup
\hypersetup{urlcolor=darkgray}
\href{\doclicenseURL}{%
\doclicenseName{} -- Attribution-NonCommercial-ShareAlike 4.0 International%
}%
\endgroup
}

% --- RESUMO (ABSTRACT) ---
\begin{abstract}
Nestas notas de aula iniciamos o estudo dos processos estocásticos a tempo
contínuo, tomando como modelo central o processo de Poisson. A construção parte
de cinco postulados sobre o número acumulado de ocorrências de um evento ao
longo do tempo, entre eles a independência das contagens em intervalos
disjuntos e a proporcionalidade da taxa de ocorrência em intervalos curtos. A
partir deles, deduzimos que o número de ocorrências até o instante $t$ tem
distribuição de Poisson de parâmetro $\lambda t$, por dois caminhos: a
resolução indutiva do sistema de equações diferenciais-diferenças e o método
das funções geradoras de probabilidade. Em seguida, estudamos os tempos de
chegada e os tempos entre chegadas. Mostramos que estes últimos são
independentes e têm distribuição exponencial. Discutimos, por fim, o papel da
propriedade de falta de memória, que caracteriza a distribuição exponencial e
permite reconstruir o processo a partir de uma sequência de exponenciais
independentes. O texto inclui figuras, exemplos, exercícios e uma lista de
problemas. A exposição segue de perto as Seções~11.1 a~11.3 do Capítulo~11 da
referência \cite{GrimmettWelsh2014}.
\end{abstract}

%----------------------------------------------------------------------
% CONTEÚDO
%----------------------------------------------------------------------

\section{Processos a Tempo Contínuo}\label{sec:tempo-continuo}

Nas notas de aula anteriores sobre processos de ramificação e sobre passeios
aleatórios, estudamos processos cujo índice era um número inteiro $n$ e que,
por isso, chamamos de processos a \emph{tempo discreto}: neles as observações
ocorrem nos instantes $0,1,2,\ldots$ e a evolução do processo é descrita por
uma sequência de variáveis aleatórias. Muitos fenômenos, no entanto, são
observados continuamente, sem que exista um intervalo mínimo entre duas
observações consecutivas. O número de chamadas telefônicas que chegam a uma
central, o número de partículas detectadas por um contador Geiger e o número de
clientes que entram em uma loja são contagens desse tipo. A partir daqui o
índice do processo deixa de ser um inteiro e passa a ser um instante $t$, e o
modelo que organiza toda a exposição é o processo de Poisson.

\subsection{Do Tempo Discreto ao Tempo Contínuo}
\label{subsec:discreto-continuo}

Um processo a tempo contínuo é uma família de variáveis aleatórias indexada pelo
conjunto dos números reais não negativos. Mais precisamente, em um espaço de
probabilidade $(\Omega,\mathcal{F},\mathbb{P})$ fixado, um processo a tempo
contínuo é uma família
\[
  Z=(Z_t)_{t\geqslant0}
\]
de variáveis aleatórias $Z_t:\Omega\to\mathbb{R}$, em que interpretamos $Z_t$
como o valor do processo no instante $t$. O caso que estudaremos é aquele em
que cada $Z_t$ assume valores inteiros, e todos os exemplos destas notas têm
essa forma.

A teoria geral dos processos a tempo contínuo é profunda, em parte porque
exige lidar com famílias não enumeráveis de variáveis aleatórias. Boa parte
dessas dificuldades desaparece quando nos restringimos a processos que assumem
apenas valores inteiros. Para o que faremos, basta que cada $Z_t$ seja uma
variável aleatória e que as trajetórias $t\mapsto Z_t(\omega)$ tenham as
propriedades de monotonia enunciadas adiante; não precisaremos construir a
família inteira por meio de um teorema de extensão.

A diferença de método em relação ao tempo discreto é simples de descrever.
Para um processo $(Z_n)_{n\geqslant0}$ a tempo discreto, relacionamos $Z_{n+1}$
com $Z_n$ e obtemos \emph{equações de recorrência} e \emph{equações de
diferenças} para as probabilidades envolvidas, como fizemos no estudo dos
passeios aleatórios. Para um processo a tempo contínuo, relacionamos $Z_{t+h}$
com $Z_t$ para um acréscimo pequeno $h>0$ e, em seguida, fazemos $h\downarrow0$.
O que era uma equação de diferenças passa a ser uma \emph{equação diferencial}.
Essa é a única mudança conceitual relevante, e ela organiza toda a exposição
seguinte.

\subsection{O Modelo da Central Telefônica}\label{subsec:modelo-central}

Apresentamos o exemplo que serve de protótipo. Bill é o chefe dos mensageiros
de um grande hotel e parte de seu trabalho consiste em atender as ligações
telefônicas que chegam. Ele não consegue prever com certeza quando o telefone
vai tocar; do seu ponto de vista, as chamadas chegam ao acaso. Fazemos duas
simplificações sobre essas chamadas. Primeiro, supomos que Bill atende cada
chamada instantaneamente, de modo que nenhuma chamada se perde, a menos que
chegue exatamente no mesmo instante que outra. Segundo, supomos que as
chamadas chegam de modo \emph{homogêneo} no tempo, no sentido de que a chance
de o telefone tocar em um período depende apenas da duração desse período. A
segunda hipótese é grosseira para um dia inteiro de trabalho, mas pode ser
razoável em certos períodos do dia.

Descrevemos o tempo por um parâmetro $t$ que percorre $[0,\infty)$ e denotamos
por $N_t$ o número de chamadas que chegaram ao hotel no intervalo $[0,t]$, isto
é, o número de chamadas que Bill atendeu até o instante $t$, inclusive.
Suponhamos que o processo $N=(N_t)_{t\geqslant0}$ evolua de modo que as
condições a seguir sejam válidas:

\begin{enumerate}[label=\Alph*., itemsep=2pt]
  \item $N_t$ é uma variável aleatória que assume valores em
        $\{0,1,2,\ldots\}$;
  \item $N_0=0$;
  \item $N_s\leqslant N_t$ sempre que $s\leqslant t$;
  \item \emph{independência}: se $0\leqslant s<t$, então o número de chamadas
        que chegam no intervalo $(s,t]$ é independente das chegadas anteriores
        ao instante $s$;
  \item \emph{taxa de chegada}: existe um número $\lambda>0$, chamado de
        \emph{taxa de chegada}, tal que, para $h>0$ pequeno,
        \begin{gather*}
          \mathbb{P}(N_{t+h}=n+1\mid N_t=n)=\lambda h+o(h),
          \\[0.2cm]
          \mathbb{P}(N_{t+h}=n\mid N_t=n)=1-\lambda h+o(h).
        \end{gather*}
\end{enumerate}

Os postulados A, B e C dizem apenas que $N_t$ conta ocorrências: o número de
chamadas é um inteiro não negativo, a contagem começa em zero e não decresce
com o tempo. O postulado D é a hipótese de independência entre o que ocorre em
um intervalo e o que ocorreu antes dele. O postulado E é a hipótese de
homogeneidade quantificada: em um intervalo curto de duração $h$, a
probabilidade de uma chamada é aproximadamente proporcional a $h$, e a
aproximação melhora quando $h$ diminui.

\subsection{O Postulado de Taxa e a Notação \texorpdfstring{$o(h)$}{o(h)}}
\label{subsec:taxa-o-h}

Antes de usar o postulado E, precisamos fixar o significado da expressão
$o(h)$. Escrevemos $f(h)=o(h)$ quando
\[
  \lim_{h\downarrow0}\frac{f(h)}{h}=0,
\]
isto é, quando $f$ é de ordem de grandeza menor que $h$ para $h$ pequeno.
Analogamente, escrevemos $f(h)=O(h)$ quando existe uma constante $C>0$ tal que
$|f(h)|\leqslant Ch$ para todo $h$ suficientemente pequeno. Nessa
notação, $o(h)$ representa uma função de $h$ que pode ser diferente a cada
aparição; em particular, a soma de duas parcelas $o(h)$ é ainda $o(h)$. Assim,
ao escrevermos
\[
  \mathbb{P}(N_{t+h}=n+1\mid N_t=n)=\lambda h+o(h),
\]
afirmamos que a diferença entre o lado esquerdo e $\lambda h$ é uma função de
$h$ cuja razão por $h$ tende a zero quando $h\downarrow0$.

O postulado E determina a probabilidade de uma chegada em um intervalo curto.
A partir dele obtemos também a probabilidade de duas ou mais chegadas no mesmo
intervalo. Como os eventos $\{N_{t+h}=n\}$, $\{N_{t+h}=n+1\}$ e
$\{N_{t+h}\geqslant n+2\}$ são dois a dois disjuntos e cobrem $\{N_t=n\}$, temos
\begin{align*}
  \mathbb{P}(N_{t+h}\geqslant n+2\mid N_t=n)
  &=1-\mathbb{P}(N_{t+h}=n\mid N_t=n)
  \\[0.2cm]
  &\quad-\mathbb{P}(N_{t+h}=n+1\mid N_t=n)
  \\[0.3cm]
  &=1-\bigl[1-\lambda h+o(h)\bigr]-\bigl[\lambda h+o(h)\bigr]
  \\[0.3cm]
  &=o(h).
\end{align*}
Concluímos que, em um intervalo curto, os únicos eventos cuja probabilidade
não é $o(h)$ são o de nenhuma chegada e o de exatamente uma chegada. A
igualdade anterior também mostra que as duas identidades do postulado E são
compatíveis entre si: somando as probabilidades dos três eventos, os dois
termos $o(h)$ se cancelam e resta $1$.

Esse é o nosso modelo para a chegada de chamadas telefônicas. Trata-se de um
modelo primitivo, baseado na ideia de chegadas ao acaso e obtido com o auxílio
de hipóteses simplificadoras. Por uma razão que ficará clara na
\Cref{sec:distribuicao-nt}, o processo $N=(N_t)_{t\geqslant0}$ é chamado de
\emph{processo de Poisson de taxa} $\lambda$. Processos de Poisson descrevem
bem muitos fenômenos, como
\begin{enumerate}[label=\textnormal{(\alph*)}, itemsep=1pt]
  \item a chegada de clientes em uma loja;
  \item os cliques emitidos por um contador Geiger ao registrar a detecção de
        partículas radioativas;
  \item a incidência de mortes em uma cidade pequena com população
        razoavelmente estável, desprezadas as variações sazonais.
\end{enumerate}
No caso da emissão de partículas radioativas, o modelo é especialmente
adequado quando a fonte tem meia-vida longa e decai lentamente.

\subsection{Trajetórias e Tempos de Chegada}\label{subsec:trajetorias}

Podemos representar os resultados de um processo de Poisson $N$ pelo gráfico de
$N_t$ em função de $t$. A \Cref{fig:trajetoria-poisson} mostra uma trajetória
típica: a função é constante por partes, cresce em saltos de uma unidade e
permanece em cada nível por um intervalo de tempo aleatório.

\begin{figure}[H]
  \centering
  \begin{tikzpicture}[x=1.32cm, y=0.98cm, >=Stealth, font=\small]

    % Eixos
    \draw[->, black!65, line width=0.55pt] (0,0) -- (10.6,0)
      node[right=2pt, font=\small, text=text] {$t$};
    \draw[->, black!65, line width=0.55pt] (0,0) -- (0,6.4)
      node[above=2pt, font=\small, text=text] {$N_t$};

    % Escada (contínua à direita): patamares [T_i, T_{i+1}) no nível i
    \draw[teal, line width=1.3pt]
      (0,0) -- (1.3,0) -- (1.3,1) -- (2.1,1) -- (2.1,2) -- (3.6,2)
      -- (3.6,3) -- (4.8,3) -- (4.8,4) -- (7.4,4) -- (7.4,5) -- (9.6,5);

    % Extremos incluídos (valor assumido no instante da chegada)
    \foreach \x/\y in {0/0, 1.3/1, 2.1/2, 3.6/3, 4.8/4, 7.4/5} {
      \fill[teal] (\x,\y) circle (1.7pt);
    }
    % Extremos excluídos (valor imediatamente antes da chegada)
    \foreach \x/\y in {1.3/0, 2.1/1, 3.6/2, 4.8/3, 7.4/4} {
      \draw[teal, fill=white, line width=0.8pt] (\x,\y) circle (1.7pt);
    }

    % Indicação de que a trajetória continua
    \draw[teal, densely dotted, line width=1.1pt] (9.6,5) -- (9.95,5);
    \node[font=\scriptsize, text=muted] at (9.78,5.32) {$\cdots$};

    % Marcas e rótulos dos instantes de chegada
    \foreach \x/\lab in {1.3/{T_1}, 2.1/{T_2}, 3.6/{T_3}, 4.8/{T_4},
                        7.4/{T_5}} {
      \draw[black!55, line width=0.45pt] (\x,-0.12) -- (\x,0.12);
      \node[below=2pt, font=\scriptsize, text=text] at (\x,-0.12) {$\lab$};
    }
    \node[below left=2pt, font=\scriptsize, text=muted] at (0,0) {$0$};

    % Níveis no eixo vertical
    \foreach \y in {1,2,3,4,5} {
      \draw[black!55, line width=0.45pt] (-0.10,\y) -- (0.10,\y);
      \node[left=2pt, font=\scriptsize, text=muted] at (-0.10,\y) {$\y$};
    }

    % Intervalo de nível constante
    \draw[teal!85!black, densely dashed, line width=0.7pt]
      (2.1,2.18) -- (3.6,2.18);
    \node[font=\scriptsize, text=teal!85!black, fill=white, inner sep=1.5pt]
      at (2.85,2.18) {$N_t=2$};

  \end{tikzpicture}
  \caption{Trajetória de um processo de Poisson. A função $N_t$ é constante
    por partes e cresce em saltos de uma unidade nos instantes aleatórios de
    chegada $T_1,T_2,\ldots$. Em cada patamar vale $N_t=i$ para
    $t\in[T_i,T_{i+1})$: o círculo fechado marca o valor assumido no instante
    da chegada e o círculo aberto marca o valor imediatamente anterior a ela.
    O segmento tracejado destaca um intervalo em que o nível se mantém igual a
    $2$, e a linha pontilhada indica que a trajetória continua.}
  \label{fig:trajetoria-poisson}
\end{figure}

Seja $T_i$ o instante em que chega a $i$-ésima chamada, de modo que
\begin{equation}\label{eq:tempo-chegada}
  T_i=\inf\{t:N_t=i\}.
\end{equation}
Temos então $0=T_0\leqslant T_1\leqslant T_2\leqslant\cdots$ e
$N_t=i$ quando $t$ pertence ao intervalo $[T_i,T_{i+1})$. A sequência
$T_0,T_1,T_2,\ldots$ é formada por variáveis aleatórias cujos valores
determinam completamente o processo: se conhecemos os $T_i$, então $N_t$ é dado
por
\[
  N_t=\max\{n:T_n\leqslant t\}.
\]

Os postulados garantem também que, com probabilidade $1$, as trajetórias
crescem em saltos de uma unidade, de modo que $T_i$ está bem definido para
todo $i$. De fato, para $h>0$, divida $[0,n]$ em intervalos de comprimento
$h$. Ao todo são $\lceil n/h\rceil$ intervalos e, em cada um deles, a
probabilidade de duas ou mais chegadas é $o(h)$; pela subaditividade, a
probabilidade de que isso ocorra em algum intervalo é, no máximo,
$\lceil n/h\rceil o(h)$, que é assintoticamente igual a $n\,o(h)/h$ e tende a
zero quando $h\downarrow0$.
Logo, com probabilidade $1$, não há duas chegadas simultâneas em $[0,n]$ e,
como $n$ é arbitrário, o mesmo vale em $[0,\infty)$.

Por isso, a sequência dos $T_i$ pode ser pensada como o \emph{processo inverso}
de $N$. A passagem de $N$ para os $T_i$ e de volta para $N$ será usada com
frequência nas próximas seções.

Duas perguntas organizam o que vem a seguir, e ambas foram enunciadas de forma
vaga até aqui. Qual é a função de probabilidade de $N_t$ para um valor fixado
de $t$? O que pode ser dito sobre a distribuição da sequência
$T_0,T_1,\ldots$? A \Cref{sec:distribuicao-nt} responde à primeira e a
\Cref{sec:tempos-chegada} responde à segunda.

\begin{exercise}\label{exercise:zero-chegadas}
Seja $N$ um processo de Poisson de taxa $\lambda$. Mostre que
\[
  \mathbb{P}(N_{t+h}=0)=\bigl[1-\lambda h+o(h)\bigr]\mathbb{P}(N_t=0)
\]
para $h>0$ pequeno. Deduza que $p_0(t)=\mathbb{P}(N_t=0)$ satisfaz a equação
diferencial
\[
  p_0'(t)=-\lambda p_0(t)
\]
e resolva essa equação para obter $p_0(t)$.
\end{exercise}

\begin{exercise}[Afinamento]\label{exercise:afinamento}
Suponha que chamadas telefônicas cheguem a uma central segundo um processo de
Poisson $N=(N_t)_{t\geqslant0}$ de taxa $\lambda$ e que o equipamento esteja
com defeito, de modo que cada chamada que chega não seja registrada com
probabilidade $q$, independentemente das demais. Seja $N'_t$ o número de
chamadas registradas até o instante $t$. Mostre que $N'=(N'_t)_{t\geqslant0}$
é um processo de Poisson de taxa $\lambda(1-q)$.
\end{exercise}

\begin{exercise}[Superposição]\label{exercise:superposicao}
Dois fluxos independentes de chamadas telefônicas chegam a uma central: o
primeiro é um processo de Poisson de taxa $\lambda$ e o segundo é um processo
de Poisson de taxa $\mu$. Mostre que o fluxo combinado de chamadas é um
processo de Poisson de taxa $\lambda+\mu$.
\end{exercise}

\section{A Distribuição do Número de Chegadas}\label{sec:distribuicao-nt}

Nesta seção calculamos a distribuição de $N_t$ para cada $t$ fixado. O
resultado justifica o nome do processo: o parâmetro da distribuição de Poisson
é $\lambda t$, proporcional ao tempo transcorrido.

\begin{theorem}\label{thm:distribuicao-nt}
Seja $N$ um processo de Poisson de taxa $\lambda$. Para cada $t>0$, a variável
aleatória $N_t$ tem distribuição de Poisson de parâmetro $\lambda t$. Isto é,
para $t>0$ e $k=0,1,2,\ldots$,
\begin{equation}\label{eq:distribuicao-nt}
  \mathbb{P}(N_t=k)=\frac{1}{k!}(\lambda t)^k e^{-\lambda t}.
\end{equation}
\end{theorem}

\begin{corollary}\label{cor:media-variancia-nt}
Para $t>0$,
\begin{equation}\label{eq:media-variancia-nt}
  \mathbb{E}[N_t]=\lambda t
  \quad\text{e}\quad
  \operatorname{Var}(N_t)=\lambda t.
\end{equation}
Em particular, a média e a variância de $N_t$ crescem linearmente com $t$, e o
quociente $\mathbb{E}[N_t]/t$ é constante e igual a $\lambda$.
\end{corollary}
\begin{proof}
Já estabelecemos, nas notas de aula sobre esperança de variáveis aleatórias
discretas e sobre funções geradoras de probabilidade, que uma variável com
distribuição $\mathrm{Poisson}(\mu)$ tem média e variância iguais a $\mu$. Por
\eqref{eq:distribuicao-nt}, $N_t$ tem essa distribuição com $\mu=\lambda t$, e
o resultado segue.
\end{proof}

A demonstração do \Cref{thm:distribuicao-nt} segue o mesmo programa que
empregamos no tempo discreto, com uma troca de ferramenta. Lá, relacionávamos
$p_k(n+1)$ com $p_k(n)$ e obtínhamos equações de diferenças; aqui,
relacionamos $p_k(t+h)$ com $p_k(t)$ e $p_{k-1}(t)$ para $h$ pequeno e
passamos ao limite $h\downarrow0$, obtendo equações diferenciais. O ponto de
partida é o Teorema da Partição, aplicado à partição do espaço amostral
determinada pelos valores de $N_t$.

\begin{proof}[Prova do \Cref{thm:distribuicao-nt}]
Escrevemos
\[
  p_k(t)=\mathbb{P}(N_t=k).
\]
Fixamos $t\geqslant0$ e tomamos $h>0$ pequeno. Para $k\geqslant1$, a
partição do espaço amostral pelos valores de $N_t$ dá
\[
  \mathbb{P}(N_{t+h}=k)
  =\sum_{i=0}^{\infty}\mathbb{P}(N_{t+h}=k\mid N_t=i)\,\mathbb{P}(N_t=i).
\]
A soma é, na verdade, finita. Pelo postulado C, se $i>k$, então
$\{N_t=i\}\cap\{N_{t+h}=k\}=\varnothing$, pois a contagem não pode decrescer.
Além disso, se $i\leqslant k-2$, então a inclusão
$\{N_{t+h}=k\}\subseteq\{N_{t+h}\geqslant i+2\}$ e a estimativa obtida na
\Cref{subsec:taxa-o-h} dão
\[
  \mathbb{P}(N_{t+h}=k\mid N_t=i)
  \leqslant\mathbb{P}(N_{t+h}\geqslant i+2\mid N_t=i)=o(h).
\]
Como o número de parcelas com $i\leqslant k-2$ é finito e não depende de $h$,
a contribuição dessas parcelas para a soma é $o(h)$. Restam, portanto, apenas
as parcelas $i=k-1$ e $i=k$, e o postulado E fornece

\begin{align}
  \mathbb{P}(N_{t+h}=k)
  &=\mathbb{P}(N_{t+h}=k\mid N_t=k-1)\,\mathbb{P}(N_t=k-1)
    \notag
  \\[0.2cm]
  &\quad+\mathbb{P}(N_{t+h}=k\mid N_t=k)\,\mathbb{P}(N_t=k)+o(h)
    \notag
  \\[0.2cm]
  &=\bigl[\lambda h+o(h)\bigr]\mathbb{P}(N_t=k-1)
    +\bigl[1-\lambda h+o(h)\bigr]\mathbb{P}(N_t=k)+o(h)
    \notag
  \\[0.2cm]
  &=\lambda h\,\mathbb{P}(N_t=k-1)
    +(1-\lambda h)\,\mathbb{P}(N_t=k)+o(h),
    \label{eq:sistema-recursivo}
\end{align}
onde, na última passagem, usamos que $o(h)\,\mathbb{P}(N_t=i)$ é $o(h)$ e que
a soma de um número finito de parcelas $o(h)$ é $o(h)$. Subtraindo
$\mathbb{P}(N_t=k)$ dos dois lados, obtemos
\[
  p_k(t+h)-p_k(t)=\lambda h\bigl[p_{k-1}(t)-p_k(t)\bigr]+o(h),
\]
para $k=1,2,\ldots$. Dividindo por $h$ e fazendo $h\downarrow0$, concluímos
que $p_k$ é derivável e que\footnote{A passagem ao limite fornece, a rigor, a
derivada à direita de $p_k$ no ponto $t$. Aplicando o mesmo argumento ao
intervalo $[t-h,t]$, isto é, tomando $t-h$ como ponto base, obtemos a mesma
identidade para a derivada à esquerda, desde que $p_{k-1}$ e $p_k$ sejam
contínuas em $t$. A continuidade, por sua vez, decorre da igualdade anterior:
existe uma constante $C$, que não depende de $h$, tal que
$\bigl|p_k(t+h)-p_k(t)\bigr|\leqslant Ch+o(h)$, e o lado direito tende a zero
quando $h\downarrow0$. Logo a igualdade vale para a derivada usual.}
\begin{equation}\label{eq:edo-nt}
  p_k'(t)=\lambda p_{k-1}(t)-\lambda p_k(t),
  \qquad k=1,2,\ldots
\end{equation}
No caso $k=0$, o mesmo argumento é mais simples, pois o evento
$\{N_{t+h}=0\}$ está contido em $\{N_t=0\}$ pelo postulado C. Assim,
\[
  \mathbb{P}(N_{t+h}=0)
  =\mathbb{P}(N_{t+h}=0\mid N_t=0)\,\mathbb{P}(N_t=0)
  =\bigl[1-\lambda h+o(h)\bigr]\mathbb{P}(N_t=0),
\]
e, repetindo o mesmo argumento (subtrair $\mathbb{P}(N_t=0)$ e dividir por
$h$),
\begin{equation}\label{eq:edo-nt-zero}
  p_0'(t)=-\lambda p_0(t).
\end{equation}
As equações \eqref{eq:edo-nt} e \eqref{eq:edo-nt-zero} formam um sistema de
equações \emph{diferenciais-diferenças} para as funções
$p_0(t),p_1(t),\ldots$, que devemos resolver com a condição inicial
decorrente do postulado B, a saber,
\begin{equation}\label{eq:condicao-inicial-nt}
  p_k(0)=
  \begin{cases}
    1, & \text{se } k=0,
    \\[0.2cm]
    0, & \text{se } k\neq0.
  \end{cases}
\end{equation}
Apresentamos duas maneiras de resolver esse sistema.

\medskip
\noindent\textbf{Solução A (indução).}
A equação \eqref{eq:edo-nt-zero} envolve apenas $p_0$. Sua solução geral é
$p_0(t)=Ae^{-\lambda t}$, e a condição \eqref{eq:condicao-inicial-nt} dá
$A=1$. Logo,
\[
  p_0(t)=e^{-\lambda t}\qquad\text{para } t\geqslant0.
\]
Substituindo em \eqref{eq:edo-nt} com $k=1$, obtemos
\[
  p_1'(t)+\lambda p_1(t)=\lambda e^{-\lambda t}.
\]
Multiplicando a equação por $e^{\lambda t}$, obtemos
$\bigl(e^{\lambda t}p_1(t)\bigr)'=\lambda$. Como $p_1(0)=0$, a integração em
$[0,t]$ dá $e^{\lambda t}p_1(t)=\lambda t$, isto é,
\[
  p_1(t)=\lambda t\,e^{-\lambda t}\qquad\text{para } t\geqslant0.
\]
Com $k=2$, o mesmo cálculo fornece
\[
  \bigl(e^{\lambda t}p_2(t)\bigr)'=\lambda\,e^{\lambda t}p_1(t)=\lambda^2t,
\]
e, integrando com $p_2(0)=0$, vem
$p_2(t)=\tfrac{1}{2}(\lambda t)^2e^{-\lambda t}$.

Os três primeiros valores sugerem a fórmula \eqref{eq:distribuicao-nt}, que
agora provamos por indução em $k$. Suponha que a fórmula valha para $k-1$,
isto é,
\[
  p_{k-1}(t)=\frac{1}{(k-1)!}(\lambda t)^{k-1}e^{-\lambda t}.
\]
Por \eqref{eq:edo-nt},
\[
  \bigl(e^{\lambda t}p_k(t)\bigr)'
  =\lambda\,e^{\lambda t}p_{k-1}(t)
  =\frac{\lambda^k}{(k-1)!}\,t^{k-1}.
\]
Como $p_k(0)=0$, a integração em $[0,t]$ dá
\[
  e^{\lambda t}p_k(t)
  =\int_0^t\frac{\lambda^k}{(k-1)!}\,s^{k-1}\,ds
  =\frac{(\lambda t)^k}{k!},
\]
e a fórmula \eqref{eq:distribuicao-nt} segue dividindo por $e^{\lambda t}$.

\medskip
\noindent\textbf{Solução B (funções geradoras).}
O segundo método é mais econômico e se aplica, com pequenas adaptações, aos
processos das próximas notas de aula. Ele usa a função geradora de
probabilidade de $N_t$, isto é,
\[
  G(s,t)=\mathbb{E}\bigl[s^{N_t}\bigr]=\sum_{k=0}^{\infty}p_k(t)\,s^k,
  \qquad |s|\leqslant1.
\]
A ideia é trocar o sistema infinito \eqref{eq:edo-nt} e \eqref{eq:edo-nt-zero}
por uma única equação diferencial para $G$. Para isso, multiplicamos
\eqref{eq:edo-nt} por $s^k$ e somamos em $k\geqslant1$. Obtemos
\[
  \sum_{k=1}^{\infty}p_k'(t)\,s^k
  =\lambda\sum_{k=1}^{\infty}p_{k-1}(t)\,s^k
   -\lambda\sum_{k=1}^{\infty}p_k(t)\,s^k.
\]
A primeira soma do lado direito é $s\,G(s,t)$, pela reindexação $j=k-1$. Para
a soma do lado esquerdo, observe que $p_0'=-\lambda p_0$ por
\eqref{eq:edo-nt-zero} e que
\[
  \sum_{k=1}^{\infty}p_k(t)\,s^k=G(s,t)-p_0(t).
\]
Somando $p_0'(t)$ aos dois lados da igualdade anterior, chegamos a
\[
  \frac{\partial G}{\partial t}(s,t)
  =\lambda s\,G(s,t)-\lambda\bigl[G(s,t)-p_0(t)\bigr]-\lambda p_0(t)
  =\lambda(s-1)\,G(s,t).
\]
A derivação termo a termo é legítima: por \eqref{eq:edo-nt},
$|p_k'(t)|\leqslant\lambda\bigl(p_{k-1}(t)+p_k(t)\bigr)$, de modo que, para
$|s|\leqslant1$,
\[
  \sum_{k=1}^{\infty}\bigl|p_k'(t)\bigr|\,|s|^k
  \leqslant2\lambda\sum_{k=0}^{\infty}p_k(t)=2\lambda,
\]
e a cota não depende de $t$. O teste de Weierstrass, aplicado às caudas dessa
série, mostra então que ela converge uniformemente em $t$: o supremo em $t$ do
módulo de cada termo é somável, com soma no máximo $2\lambda$. Para cada $s$
fixado com $|s|\leqslant1$, a função $t\mapsto G(s,t)$ é derivável e sua
derivada é a série das derivadas.

Resta resolver a equação diferencial ordinária
\[
  \frac{\partial G}{\partial t}=\lambda(s-1)\,G,
\]
na variável $t$, com $s$ fixado. Ela é linear de primeira ordem e se resolve
por fator integrante. Multiplicando por $e^{-\lambda t(s-1)}$, obtemos
\[
  \frac{\partial}{\partial t}\Bigl(e^{-\lambda t(s-1)}G(s,t)\Bigr)=0,
\]
de modo que $e^{-\lambda t(s-1)}G(s,t)=A(s)$ depende apenas de $s$. A condição
inicial vem de \eqref{eq:condicao-inicial-nt}:
\[
  G(s,0)=\sum_{k=0}^{\infty}p_k(0)\,s^k=1,
\]
para todo $s$, logo $A(s)=1$. Portanto
\[
  G(s,t)=e^{\lambda t(s-1)}
  =e^{-\lambda t}\sum_{k=0}^{\infty}\frac{(\lambda t)^k}{k!}\,s^k,
\]
onde usamos a série da exponencial. O lado direito é uma série de potências em
$s$ e, pela unicidade dos coeficientes de uma série de potências, o coeficiente
de $s^k$ coincide com $p_k(t)$. Assim,
\[
  p_k(t)=\frac{1}{k!}(\lambda t)^k e^{-\lambda t},
\]
como queríamos.
\end{proof}

A fórmula \eqref{eq:distribuicao-nt} mostra que $N_t$ tem distribuição
$\mathrm{Poisson}(\lambda t)$: o parâmetro cresce com o tempo, e não é fixo. A
mesma família de distribuições aparece nas notas de aula sobre funções
geradoras de probabilidade, onde a função geradora $e^{\mu(s-1)}$ foi
identificada como a da distribuição de Poisson de parâmetro $\mu$; aqui
$\mu=\lambda t$. A \Cref{fig:poisson-hastes} compara as funções de
probabilidade de $N_1$ e de $N_4$ no caso $\lambda=1$.

\begin{figure}[htbp]
  \centering
  \begin{tikzpicture}[x=1.15cm, y=10cm, >=Stealth, font=\small]

    % Eixos
    \draw[->, black!65, line width=0.55pt] (-0.80,0) -- (9.95,0)
      node[right=2pt, font=\small, text=text] {$k$};
    \draw[->, black!65, line width=0.55pt] (-0.80,0) -- (-0.80,0.45)
      node[above=2pt, font=\small, text=text] {$\mathbb{P}(N_t=k)$};

    % Escalas
    \foreach \k in {0,...,9} {
      \draw[black!55, line width=0.45pt] (\k,-0.008) -- (\k,0.008);
      \node[below=2pt, font=\scriptsize, text=muted] at (\k,-0.008) {$\k$};
    }
    \foreach \v/\lab in {0.1/{0{,}1}, 0.2/{0{,}2}, 0.3/{0{,}3}} {
      \draw[black!55, line width=0.45pt] (-0.88,\v) -- (-0.72,\v);
      \node[left=2pt, font=\scriptsize, text=muted] at (-0.88,\v) {$\lab$};
    }

    % t = 1, lambda t = 1
    \foreach \k/\h in {0/0.3679, 1/0.3679, 2/0.1839, 3/0.0613, 4/0.0153,
                      5/0.0031} {
      \filldraw[fill=teal!65, draw=teal, rounded corners=0.6pt]
        ({\k-0.31},0) rectangle ({\k-0.03},\h);
    }

    % t = 4, lambda t = 4
    \foreach \k/\h in {0/0.0183, 1/0.0733, 2/0.1465, 3/0.1954, 4/0.1954,
                      5/0.1563, 6/0.1042, 7/0.0595, 8/0.0298, 9/0.0132} {
      \filldraw[fill=muted!22, draw=muted!70, rounded corners=0.6pt]
        ({\k+0.03},0) rectangle ({\k+0.31},\h);
    }

    % Legenda
    \filldraw[fill=teal!65, draw=teal] (6.55,0.405) rectangle (6.83,0.437);
    \node[right=3pt, font=\scriptsize, text=text] at (6.83,0.421)
      {$\lambda t=1$};
    \filldraw[fill=muted!22, draw=muted!70] (6.55,0.365) rectangle (6.83,0.397);
    \node[right=3pt, font=\scriptsize, text=text] at (6.83,0.381)
      {$\lambda t=4$};

  \end{tikzpicture}
  \caption{Funções de probabilidade de $N_t$ para $\lambda=1$, com $t=1$
    (barras escuras) e $t=4$ (barras cinzas). Ao crescer $t$, as barras se
    deslocam para a direita e a distribuição se espalha, acompanhando o
    crescimento da média e da variância, ambas iguais a $\lambda t$.}
  \label{fig:poisson-hastes}
\end{figure}

\begin{remark}\label{rem:tempos-finitos}
Como consequência de \eqref{eq:distribuicao-nt}, para cada $k$ fixado,
\[
  \mathbb{P}(N_t\leqslant k)
  =e^{-\lambda t}\sum_{j=0}^{k}\frac{(\lambda t)^j}{j!}
  \longrightarrow0
  \qquad\text{quando } t\to\infty.
\]
Como a trajetória $t\mapsto N_t$ é não decrescente,
\[
  \Bigl\{\sup_{t\geqslant0}N_t\leqslant k\Bigr\}
  =\bigcap_{n=1}^{\infty}\{N_n\leqslant k\},
\]
e os eventos do lado direito formam uma sequência decrescente. Pela
continuidade da probabilidade e por \eqref{eq:distribuicao-nt},
\[
  \mathbb{P}\Bigl(\sup_{t\geqslant0}N_t\leqslant k\Bigr)
  =\lim_{n\to\infty}\mathbb{P}(N_n\leqslant k)=0.
\]
A união desses eventos sobre todos os $k$ também tem probabilidade nula, logo
$N_t\to\infty$ com probabilidade $1$ quando $t\to\infty$. Em particular, cada
tempo de chegada $T_i$ definido em \eqref{eq:tempo-chegada} é finito quase
certamente, o que legitima o uso da sequência $T_0,T_1,\ldots$ na
\Cref{sec:tempos-chegada}.
\end{remark}

\begin{exercise}\label{exercise:variancia-quociente}
Seja $N$ um processo de Poisson de taxa $\lambda$. Mostre que
$\operatorname{Var}(N_t/t)\to0$ quando $t\to\infty$.
\end{exercise}

A função geradora $G(s,t)$ também permite calcular probabilidades que envolvem
$N_t$ sem passar pela fórmula fechada, quando $G$ é avaliada em valores
especiais de $s$. O exercício a seguir é um exemplo: para obter a
probabilidade de $N_t$ ser par, avalie $G(s,t)$ em $s=-1$ e use que os eventos
$\{N_t\text{ é par}\}$ e $\{N_t\text{ é ímpar}\}$ são complementares.

\begin{exercise}\label{exercise:paridade-poisson}
Seja $N$ um processo de Poisson de taxa $\lambda$. Mostre que, para $t>0$,
\[
  \mathbb{P}(N_t\text{ é par})=e^{-\lambda t}\cosh(\lambda t)
  \qquad\text{e}\qquad
  \mathbb{P}(N_t\text{ é ímpar})=e^{-\lambda t}\operatorname{senh}(\lambda t).
\]
\end{exercise}

\begin{exercise}\label{exercise:fgm-clt-poisson}
Seja $N$ um processo de Poisson de taxa $\lambda$ e considere a variável
aleatória padronizada
\[
  U_t=\frac{N_t-\mathbb{E}[N_t]}{\sqrt{\operatorname{Var}(N_t)}}.
\]
Mostre que a função geradora de momentos de $U_t$ é
\[
  M_t(x)=\mathbb{E}\bigl[e^{xU_t}\bigr]
  =\exp\Bigl[-x\sqrt{\lambda t}+\lambda t\bigl(e^{x/\sqrt{\lambda t}}-1\bigr)\Bigr],
\]
para $x\in\mathbb{R}$. Deduza que, quando $t\to\infty$,
\[
  \mathbb{P}(U_t\leqslant u)
  \longrightarrow
  \int_{-\infty}^{u}\frac{1}{\sqrt{2\pi}}e^{-\frac{1}{2}v^2}\,dv
  \qquad\text{para todo } u\in\mathbb{R}.
\]
Este é o Teorema do Limite Central para o processo de Poisson.
\end{exercise}

\section{Tempos entre Chegadas e a Distribuição Exponencial}
\label{sec:tempos-chegada}

Na seção anterior descrevemos a distribuição de $N_t$ para cada instante $t$
fixado. Passamos agora à segunda pergunta anunciada na
\Cref{subsec:trajetorias}: como se distribuem os tempos de chegada
$T_0,T_1,\ldots$? Veremos que os intervalos entre chegadas sucessivas são
variáveis aleatórias independentes com distribuição exponencial, e que a
propriedade que torna essa descrição possível é a falta de memória.

\subsection{Tempos de Chegada e Tempos entre Chegadas}
\label{subsec:tempos-entre-chegadas}

Seja $N$ um processo de Poisson de taxa $\lambda$. Recordamos de
\eqref{eq:tempo-chegada} que os \emph{tempos de chegada} de $N$ são
$T_i=\inf\{t:N_t=i\}$, com $T_0=0$. Em outras palavras, $T_i$ é o instante da
chegada da $i$-ésima chamada. Os \emph{tempos entre chegadas} $X_1,X_2,\ldots$
são as durações dos intervalos entre chegadas sucessivas:
\begin{equation}\label{eq:tempo-entre-chegadas}
  X_i=T_i-T_{i-1},
  \qquad i=1,2,\ldots
\end{equation}
A \Cref{fig:tempos-entre-chegadas} ilustra a relação entre as duas sequências.
Os tempos de chegada são as somas acumuladas dos tempos entre chegadas,
\[
  T_k=X_1+X_2+\cdots+X_k,
\]
e os tempos entre chegadas são as diferenças sucessivas dos tempos de chegada.

\begin{figure}[H]
  \centering
  \begin{tikzpicture}[x=1.95cm, y=1cm, >=Stealth, font=\small]

    % Eixo dos tempos
    \draw[->, black!65, line width=0.55pt] (-0.25,0) -- (6.4,0)
      node[right=2pt, font=\small, text=text] {$t$};

    % Tempos de chegada
    \foreach \x in {0, 1.3, 2.1, 3.6, 4.8} {
      \fill[teal] (\x,0) circle (2pt);
    }
    \node[above=7pt, font=\scriptsize, text=text] at (0,0) {$T_0=0$};
    \node[above=7pt, font=\scriptsize, text=text] at (1.3,0) {$T_1$};
    \node[above=7pt, font=\scriptsize, text=text] at (2.1,0) {$T_2$};
    \node[above=7pt, font=\scriptsize, text=text] at (3.6,0) {$T_3$};
    \node[above=7pt, font=\scriptsize, text=text] at (4.8,0) {$T_4$};

    % Tempos entre chegadas
    \foreach \a/\b/\lab in {0/1.3/{X_1}, 1.3/2.1/{X_2}, 2.1/3.6/{X_3},
                            3.6/4.8/{X_4}} {
      \draw[teal!85!black, line width=0.6pt, decorate,
        decoration={brace, amplitude=3.5pt, mirror}]
        (\a,-0.28) -- (\b,-0.28);
      \node[below=9pt, font=\scriptsize, text=teal!85!black]
        at ({(\a+\b)/2},-0.28) {$\lab$};
    }

    % Indicação de que a sequência continua
    \draw[black!45, densely dotted, line width=0.6pt] (4.8,0) -- (5.7,0);
    \node[font=\scriptsize, text=muted] at (5.95,0.05) {$\cdots$};

  \end{tikzpicture}
  \caption{Relação entre os tempos de chegada e os tempos entre chegadas. Os
    pontos sobre o eixo marcam os instantes $T_0,T_1,\ldots$ em que ocorrem as
    chegadas, e as chaves indicam os intervalos $X_i=T_i-T_{i-1}$ entre
    chegadas sucessivas; a linha pontilhada indica que a sequência continua.
    Conhecer uma das duas sequências determina a outra.}
  \label{fig:tempos-entre-chegadas}
\end{figure}

\begin{theorem}\label{thm:exponenciais-iid}
Em um processo de Poisson de taxa $\lambda$, os tempos entre chegadas
$X_1,X_2,\ldots$ são variáveis aleatórias independentes, cada uma com
distribuição exponencial de parâmetro $\lambda$.
\end{theorem}

Esse resultado liga os postulados do processo de Poisson à distribuição
exponencial. Nos processos que satisfazem uma condição de independência como o
postulado D, a distribuição dos tempos de espera é exponencial, e a razão é a
propriedade de falta de memória, que discutiremos na
\Cref{subsec:falta-memoria} e que caracteriza a exponencial entre as
distribuições de variáveis positivas.

\begin{proof}[Prova do \Cref{thm:exponenciais-iid}]
Consideremos primeiro $X_1$. Para $u\geqslant0$, o evento $\{X_1>u\}$ é o
evento de não haver chegada no intervalo $[0,u]$, isto é,
\[
  \mathbb{P}(X_1>u)=\mathbb{P}(N_u=0).
\]
Pelo \Cref{thm:distribuicao-nt}, a cauda de $X_1$ é
\[
  \mathbb{P}(X_1>u)=e^{-\lambda u},
  \qquad u\geqslant0,
\]
e portanto $X_1$ tem distribuição exponencial de parâmetro $\lambda$.

Pelo postulado D, as chegadas no intervalo $(0,X_1]$ são independentes das
chegadas posteriores a $X_1$. Segue que o tempo de espera $X_2$ pela chegada
seguinte é independente de $X_1$. Além disso, as chegadas ocorrem de modo
homogêneo no tempo, pois a probabilidade de uma chegada em $(t,t+h]$ depende
apenas de $h$ e não de $t$, como registra o postulado E. Logo $X_2$ tem a
mesma distribuição que $X_1$. Repetindo o argumento, todos os $X_i$ são
independentes e têm a mesma distribuição que $X_1$.
\end{proof}

A demonstração acima está incompleta, e vale registrar por quê. O passo que
usa a independência trata de um intervalo $(0,X_1]$ de comprimento
\emph{aleatório}, e o postulado D foi enunciado para intervalos de
comprimento fixo. Tornar esse passo rigoroso não é imediato: é preciso
mostrar que o processo ``recomeça do zero'' em um instante aleatório, e essa é
a chamada \emph{propriedade de Markov forte}, que estudaremos em notas de aula
futuras, primeiramente no contexto de processos a tempo discreto. Por essa
razão, apresentamos apenas um esboço de prova.

A \Cref{prop:construcao-poisson} mostra que a recíproca vale, isto é, que
uma sequência de exponenciais independentes produz de fato um processo de
Poisson. Para os fins de modelagem, os dois pontos de vista são equivalentes.

\subsection{A Propriedade de Falta de Memória}\label{subsec:falta-memoria}

A distribuição exponencial ocupa um lugar especial entre as distribuições
contínuas por causa da seguinte propriedade, que já encontramos nas notas de
aula sobre vetores contínuos.

\begin{definition}\label{def:falta-memoria}
Dizemos que uma variável aleatória $X$ com $\mathbb{P}(X>0)=1$ tem a
propriedade de \emph{falta de memória} se
\begin{equation}\label{eq:falta-memoria}
  \mathbb{P}(X>u+v\mid X>u)=\mathbb{P}(X>v)
\end{equation}
para todos $u,v\geqslant0$ tais que $\mathbb{P}(X>u)>0$.
\end{definition}

A condição $\mathbb{P}(X>u)>0$ garante que a probabilidade condicional esteja
definida. Quando $X$ tem suporte limitado, ela pode falhar para $u$ grande; e,
nesse caso, nada há a exigir, pois o evento $\{X>u\}$ tem probabilidade nula.

Pensando em $X$ como o tempo que transcorre até a ocorrência de um evento $A$,
a condição \eqref{eq:falta-memoria} exige que, se $A$ ainda não ocorreu no
instante $u$, então o tempo adicional até a ocorrência de $A$ não dependa de
$u$. Em outras palavras, a distribuição do tempo adicional até a ocorrência de
$A$ não depende de $u$: $X$ não se lembra de quanto tempo já passou.

\begin{theorem}\label{thm:falta-memoria-caracteriza}
Seja $X$ uma variável aleatória com $\mathbb{P}(X>0)=1$. Então $X$ tem a
propriedade de falta de memória se, e somente se, existe $\lambda>0$ tal que
$X\sim\mathrm{Exp}(\lambda)$.
\end{theorem}

\begin{proof}
Suponhamos primeiro que $X\sim\mathrm{Exp}(\lambda)$. Para $u,v\geqslant0$, a
definição de probabilidade condicional, a inclusão
$\{X>u+v\}\subseteq\{X>u\}$ e a forma $e^{-\lambda x}$ da cauda dão
\begin{align*}
  \mathbb{P}(X>u+v\mid X>u)
  &=\frac{\mathbb{P}(X>u+v)}{\mathbb{P}(X>u)}
  \\[0.3cm]
  &=\frac{e^{-\lambda(u+v)}}{e^{-\lambda u}}
  =e^{-\lambda v}
  =\mathbb{P}(X>v),
\end{align*}
e $X$ tem a propriedade de falta de memória.

Reciprocamente, suponha que $X$ tenha a propriedade e escreva
$G(u)=\mathbb{P}(X>u)$. Para $u$ com $G(u)>0$, o lado esquerdo de
\eqref{eq:falta-memoria} é $G(u+v)/G(u)$, de modo que $G$ satisfaz a equação
funcional
\begin{equation}\label{eq:equacao-funcional-falta-memoria}
  G(u+v)=G(u)\,G(v)
\end{equation}
sempre que $G(u)>0$. Como $G$ é não crescente, $G(0)=1$ e $G(n)\to0$ quando
$n\to\infty$, as possibilidades $G(1)=0$ e $G(1)=1$ são incompatíveis com a
equação funcional, como mostramos nas notas citadas. Logo
$0<G(1)<1$. Tomando $\lambda=-\log G(1)>0$, a equação funcional dá
$G(q)=e^{-\lambda q}$ para todo racional não negativo $q$, e a monotonicidade
estende a identidade a todo $u\geqslant0$. Portanto $X$ tem distribuição
exponencial de parâmetro $\lambda$.
\end{proof}

A falta de memória é o que permite calcular probabilidades que envolvem
tempos de espera sem conhecer o instante em que a espera começou. No
contexto do processo de Poisson, ela diz que, se já esperamos $u$ unidades de
tempo sem chegada, a espera adicional tem a mesma distribuição que a espera
inicial. Essa é a propriedade que torna consistente a descrição do processo
pelos postulados A a E: graças a ela, o processo pode recomeçar, a partir de
qualquer instante fixado, com a mesma lei que tinha no instante inicial.

\subsection{A Recíproca: Construindo o Processo a Partir das
Exponenciais}\label{subsec:construcao-exponenciais}

Vimos que os tempos entre chegadas de um processo de Poisson são
independentes e exponenciais. Vale a recíproca: uma sequência de exponenciais
independentes, acumulada, produz um processo de Poisson. Esse é o sentido em
que a distribuição exponencial caracteriza o processo.

\begin{proposition}\label{prop:construcao-poisson}
Sejam $X_1^*,X_2^*,\ldots$ variáveis aleatórias independentes, cada uma com
distribuição $\mathrm{Exp}(\lambda)$, com $\lambda>0$. Defina, para
$k\geqslant1$ e $t\geqslant0$,
\begin{gather*}
  T_k^*=X_1^*+X_2^*+\cdots+X_k^*,
  \\[0.2cm]
  N_t^*=\max\{k:T_k^*\leqslant t\},
\end{gather*}
com a convenção de que o máximo do conjunto vazio é $0$. Então
$N^*=(N_t^*)_{t\geqslant0}$ é um processo de Poisson de taxa $\lambda$.
\end{proposition}

\begin{proof}
Antes de verificar os postulados, é preciso garantir que $N_t^*$ esteja bem
definido, isto é, que $T_k^*\to\infty$ quase certamente. Como
$\mathbb{E}[X_1^*]=1/\lambda<\infty$, a Lei Forte dos Grandes Números,
demonstrada nas notas de aula correspondentes, garante que
\[
  \frac{T_k^*}{k}=\frac{1}{k}\sum_{i=1}^{k}X_i^*
  \longrightarrow\frac{1}{\lambda}
  \qquad\text{quase certamente.}
\]
Em particular, $T_k^*\to\infty$ quase certamente e, para cada $t$, o conjunto
$\{k:T_k^*\leqslant t\}$ é finito, de modo que o máximo que define $N_t^*$
existe.

Verificamos agora os postulados A a E.

\smallskip
\noindent\textbf{Parte 1 (distribuição de $N_t^*$).}
Como $\mathrm{Exp}(\lambda)$ é a distribuição $\mathrm{Gama}(1,\lambda)$,
podemos usar o exemplo das notas de aula sobre somas de variáveis aleatórias
contínuas: a soma de variáveis aleatórias independentes com distribuição Gama
de mesmo parâmetro tem distribuição Gama com os parâmetros somados. Logo
$T_k^*\sim\mathrm{Gama}(k,\lambda)$. Assim, para $x>0$,
\[
  f_{T_k^*}(x)=\frac{\lambda^k}{(k-1)!}\,x^{k-1}e^{-\lambda x}.
\]
A função de distribuição correspondente é
\[
  \mathbb{P}(T_k^*\leqslant t)
  =1-e^{-\lambda t}\sum_{j=0}^{k-1}\frac{(\lambda t)^j}{j!}.
\]
De fato, as duas funções de $t$ têm a mesma derivada,
$\lambda^k t^{k-1}e^{-\lambda t}/(k-1)!$, e coincidem em $t=0$. Logo, para
$k\geqslant1$,
\begin{align*}
  \mathbb{P}(N_t^*=k)
  &=\mathbb{P}\bigl(T_k^*\leqslant t<T_{k+1}^*\bigr)
  \\
  &=\mathbb{P}(T_k^*\leqslant t)-\mathbb{P}(T_{k+1}^*\leqslant t)
  \\[0.3cm]
  &=e^{-\lambda t}\frac{(\lambda t)^k}{k!},
\end{align*}
e, para $k=0$, a mesma fórmula vale, pois $\mathbb{P}(N_t^*=0)=e^{-\lambda t}$.
Logo $N_t^*\sim\mathrm{Poisson}(\lambda t)$, e em particular valem os
postulados A e B.

\smallskip
\noindent\textbf{Parte 2 (falta de memória do tempo restante).}
Fixe $s\geqslant0$ e condicione no que ocorreu até o instante $s$, isto é, no
valor $N_s^*=k$ e no instante $T_k^*=u\leqslant s$ da última chegada. O tempo
que falta até a próxima chegada é $R_s=T_{k+1}^*-s=X_{k+1}^*-(s-u)$. Como
$X_{k+1}^*$ é independente de $X_1^*,\ldots,X_k^*$ e
$\mathbb{P}(X_{k+1}^*>s-u)>0$, a falta de memória dá
\[
  \mathbb{P}(R_s>h\mid N_s^*=k,\,T_k^*=u)
  =\mathbb{P}\bigl(X_{k+1}^*>s-u+h\mid X_{k+1}^*>s-u\bigr)
  =e^{-\lambda h}
\]
para todo $h\geqslant0$. O lado direito não depende de $k$, de $u$ nem de
$s$: condicionado ao passado, o tempo restante é $\mathrm{Exp}(\lambda)$ e é
independente desse passado. Além disso, os tempos entre chegadas
$X_{k+2}^*,X_{k+3}^*,\ldots$ são independentes do passado e do tempo
residual, de modo que toda a sequência de chegadas posterior a $s$ tem a mesma
lei e é independente da trajetória até $s$. Daí seguem os postulados D e E,
pois
\[
  \mathbb{P}(N_{s+h}^*=k\mid N_s^*=k)=e^{-\lambda h}=1-\lambda h+o(h)
\]
e, como $\{N_{s+h}^*=k+1\}\subseteq\{R_s\leqslant h\}$,
\begin{align*}
  \mathbb{P}(N_{s+h}^*=k+1\mid N_s^*=k)
  &=\mathbb{P}(R_s\leqslant h)-\mathbb{P}\bigl(R_s\leqslant h,\,
    T_{k+2}^*\leqslant s+h\bigr)
  \\[0.3cm]
  &=\bigl[\lambda h+o(h)\bigr]-o(h)
  =\lambda h+o(h),
\end{align*}
onde o segundo termo é $o(h)$ porque exige duas chegadas em $(s,s+h]$: dois
tempos restantes consecutivos menores que $h$, condicionalmente independentes,
cada um com probabilidade $O(h)$, de modo que o produto das duas
probabilidades é $O(h^2)=o(h)$.

\medskip
\noindent\textbf{Parte 3 (monotonicidade).}
Por construção, $N_t^*$ conta o número de termos da sequência $T_k^*$ que não
excedem $t$. Como $T_1^*\leqslant T_2^*\leqslant\cdots$, a contagem é não
decrescente em $t$, o que é o postulado C.
\end{proof}

A proposição anterior mostra que os postulados A a E descrevem a mesma classe
de processos que a construção por exponenciais independentes. É por isso que
podemos escolher, em cada aplicação, a descrição mais conveniente: os
postulados são úteis para deduzir equações diferenciais, e a construção por
exponenciais é útil para simular trajetórias e para calcular tempos de espera.

\begin{example}[Esperando o ônibus]\label{ex:esperando-onibus}
Suponha que os ônibus para o centro cheguem ao ponto da esquina segundo um
processo de Poisson. Sabendo disso, Davi espera um tempo exponencialmente
distribuído até que um ônibus o apanhe. Se ele chega ao ponto e Doris lhe diz
que espera o ônibus há $50$ minutos, isso não é uma boa nem uma má notícia para
Davi, pois a distribuição exponencial tem a propriedade de falta de memória.
Do mesmo modo, se ele chega bem no momento em que um ônibus parte, não precisa
temer que sua espera seja maior que o habitual. Essas propriedades são
características do processo de Poisson e não valem, por exemplo, quando os
ônibus seguem um horário fixo.
\end{example}

\begin{exercise}\label{exercise:soma-poisson-processos}
Sejam $M$ e $N$ processos de Poisson independentes, com taxas $\lambda$ e
$\mu$, respectivamente. Recorde, das notas de aula sobre independência de
variáveis aleatórias discretas, que a soma de variáveis independentes com
distribuição de Poisson tem distribuição de Poisson. Use esse resultado para
mostrar que $M+N=(M_t+N_t)_{t\geqslant0}$ é um processo de Poisson de taxa
$\lambda+\mu$. Compare esse método com o do \Cref{exercise:superposicao}.
\end{exercise}

\begin{exercise}\label{exercise:clt-via-tempos-chegada}
Se $T_k$ é o instante da $k$-ésima chegada do processo de Poisson $N$, mostre
que $N_t<k$ se, e somente se, $T_k>t$. Use o \Cref{thm:exponenciais-iid} e o
Teorema do Limite Central, demonstrado nas notas de aula correspondentes, para
deduzir que, quando $t\to\infty$,
\[
  \mathbb{P}\Bigl(\frac{N_t-\lambda t}{\sqrt{\lambda t}}\leqslant u\Bigr)
  \longrightarrow
  \int_{-\infty}^{u}\frac{1}{\sqrt{2\pi}}e^{-\frac{1}{2}v^2}\,dv
  \qquad\text{para todo } u\in\mathbb{R}.
\]
Compare o resultado com o do \Cref{exercise:fgm-clt-poisson}.
\end{exercise}

\begin{exercise}\label{exercise:chamadas-alternadas}
Chamadas chegam a uma central segundo um processo de Poisson de taxa
$\lambda$, mas a telefonista se distrai com frequência e atende apenas uma a
cada duas chamadas. Qual é a distribuição comum dos intervalos de tempo entre
chamadas sucessivas que recebem atendimento?
\end{exercise}

\section{Problemas}\label{sec:problemas}

Os problemas a seguir ampliam o que desenvolvemos nas seções anteriores. Os
quatro primeiros exploram consequências diretas dos postulados e da
distribuição de $N_t$; os dois últimos introduzem duas variantes do processo
de Poisson que aparecem com frequência em aplicações, o processo de Poisson
composto e o processo de Poisson não homogêneo. Não apresentamos soluções.
Quando necessário, indicamos em cada enunciado qual resultado já disponível
deve ser empregado.

\providecommand{\starredproblem}{%
\stepcounter{enumi}\llap{\textasteriskcentered\,}\arabic{enumi}.}

\begin{enumerate}[label={\arabic*.}]
\item\label{prob:distribuicao-incremento}
Seja $N$ um processo de Poisson de taxa $\lambda$. Qual é a distribuição de
$N_t-N_s$ para $0\leqslant s\leqslant t$? Justifique a resposta a partir dos
postulados.

\item\label{prob:covariancia-nt}
Seja $N$ um processo de Poisson de taxa $\lambda$. Mostre que
\[
  \operatorname{Cov}(N_s,N_t)=\lambda s
  \qquad\text{para } 0\leqslant s<t.
\]
Sugestão: escreva $N_t=N_s+(N_t-N_s)$ e use a independência entre os
incrementos, juntamente com \eqref{eq:media-variancia-nt}.

\item\label{prob:robos}
Três robôs aparentemente idênticos, chamados James, Simon e John, são postos a
trabalhar no instante $t=0$. Em seguida, cada um para de funcionar após um
tempo aleatório, independentemente dos outros dois, e a probabilidade de que um
dado robô pare no intervalo $(t,t+h]$, dado que ele ainda funciona no instante
$t$, é $\lambda h+o(h)$, com $\lambda>0$. Em símbolos, se $X$ é
o tempo de funcionamento, então
\[
  \mathbb{P}(t<X\leqslant t+h\mid X>t)=\lambda h+o(h).
\]
Mostre que cada robô funciona por um período de tempo com distribuição
exponencial de parâmetro $\lambda$ e que a probabilidade de ao menos um dos
três ter parado até o instante $t$ é $1-e^{-3\lambda t}$. Qual é a
probabilidade de que eles parem de funcionar na ordem James, Simon, John?

\item[\starredproblem]\label{prob:excesso-de-vida}
Sejam $X_1,X_2,\ldots$ variáveis aleatórias independentes e identicamente
distribuídas, com função de distribuição
\[
  F(x)=
  \begin{cases}
    1-e^{-\lambda x}, & \text{se } x\geqslant0,
    \\[0.2cm]
    0, & \text{se } x<0,
  \end{cases}
\]
em que $\lambda>0$. Seja $S_n=X_1+X_2+\cdots+X_n$. Prove que $S_n$ tem função
densidade
\[
  f_n(x)=\frac{1}{(n-1)!}\,\lambda^n x^{n-1}e^{-\lambda x},
  \qquad x\geqslant0.
\]
Deduza que $N_t=\max\{n:S_n\leqslant t\}$ tem distribuição de Poisson de
parâmetro $\lambda t$. O \emph{excesso de vida} $e_t$ é definido por
\[
  e_t=S_{N_t+1}-t.
\]
Se $g(t,x)=\mathbb{P}(e_t>x)$, mostre, condicionando no valor de $X_1$, que
\[
  g(t,x)=e^{-\lambda(t+x)}+\int_0^t g(t-u,x)\,\lambda e^{-\lambda u}\,du.
\]
Encontre uma solução dessa equação.

\item\label{prob:poisson-composto}
Ônibus de excursão chegam a um ponto turístico segundo um processo de Poisson
de taxa $\lambda$, e o número de turistas em cada ônibus é uma variável
aleatória com função geradora de probabilidade $G(s)$; suponha que esses
números sejam independentes entre si, identicamente distribuídos e
independentes do processo de chegadas. Mostre que o número total de turistas
que chegaram ao ponto até o instante $t$ tem função geradora de probabilidade
\[
  \exp\bigl(\lambda t\,[G(s)-1]\bigr).
\]
Esse é um exemplo do chamado processo de Poisson \emph{composto}.

\item\label{prob:poisson-nao-homogeneo}
A probabilidade de ocorrer uma falha em um sistema no intervalo de tempo
$(t,t+\tau)$ é $\lambda(t)\tau+o(\tau)$, independentemente das falhas
anteriores, e a probabilidade de ocorrer mais de uma falha nesse intervalo é
$o(\tau)$, em que $\lambda$ é uma função positiva e integrável, chamada de
\emph{função de taxa}. Prove que o número de falhas em $(0,t)$ tem
distribuição de Poisson com média $\int_0^t\lambda(x)\,dx$.

Seja $T$ o instante de ocorrência da primeira falha. Encontre a função
densidade de $T$ e mostre que, se $\lambda(t)=c/(1+t)$, com $c>0$, então
$\mathbb{E}[T]$ é finito se, e somente se, $c>1$. Esse é um exemplo do chamado
processo de Poisson \emph{não homogêneo}.
\end{enumerate}

%----------------------------------------------------------------------
% BIBLIOGRAFIA
%----------------------------------------------------------------------
\bibliographystyle{plain}
\bibliography{../arquivos-aux/referencias}

\end{document}
